%---------------------------------------------------------------------
\chapter{Agentic Concepts Demonstrated}
\label{ch:concepts}
%---------------------------------------------------------------------

\section{Concept Coverage Matrix}
Table~\ref{tab:concepts} maps each concept required by the course to the part of the system that
demonstrates it and to the evidence a reviewer can inspect. A concept that is not demonstrated is marked as such.

{\small
\begin{longtable}{@{}p{2.9cm} p{1.9cm} p{4.9cm} p{4.2cm}@{}}
\caption{Coverage of the course concepts and where the evidence is}\label{tab:concepts}\\
\toprule
\textbf{Concept} & \textbf{Tier} & \textbf{Demonstration in ReCycleGuard} & \textbf{Evidence} \\ \midrule
\endfirsthead
\toprule
\textbf{Concept} & \textbf{Tier} & \textbf{Demonstration} & \textbf{Evidence} \\ \midrule
\endhead
Tool use & Mandatory & Six whitelisted tools; the reconciler runs exact decimal arithmetic and the model reads the result & Tool intent and result records in the audit trail; tool unit tests \\
Agentic RAG & Mandatory & Graded BM25 retrieval, temporal filtering, bounded expansion, up to two reformulations and escalation & Retrieval attempts and grades in the trace; retrieval regression tests \\
Planning and reasoning & Mandatory & Explicit PLAN and ASSESS phases, a stored plan before each action, and a stored assessment after it & Plan and assessment records per turn; run states in Figure~\ref{fig:states} \\
Observability & Mandatory & Append-only event stream, exact prompt and reply per call, UI timeline, SSE stream & Audit export; trace screens (Chapter~\ref{ch:eval}) \\
Evaluation & Mandatory & Independent oracle, exact-balance scoring, finding precision and recall, abstention accuracy, human rubric & Scorer and tables (Chapter~\ref{ch:eval}) \\
Multi-agent collaboration & Bonus & Not demonstrated. A collector and auditor pair is described in the design as an optional experiment and was not built or evaluated & --- \\
Memory & Bonus & Within a run only: persisted agent state and previous rejection reasons are passed back to the model. There is no memory across runs & Checkpointed state; prompt field \code{previous\_rejection} \\
Human in the loop and guardrails & Bonus & Mandatory reviewer decision bound to digests; whitelist; bounded loops; finding gate; prompt-injection handling & Review service tests; gate rejections; budget enforcement test \\
Durable workflow & Bonus & Turn-by-turn checkpointed graph; resume after a killed process without repeating a tool call & Resume test after \code{kill -9}; Section~\ref{sec:crash} \\
Experimentation & Bonus & Controlled comparison of fixed and adaptive investigation with identical tools, corpus and budgets & Chapter~\ref{ch:exp} (held-out result pending) \\
\bottomrule
\end{longtable}}

\section{Evidence for Each Concept}

\subsection{Tool Use}
Each tool has a typed request, a handler and a result that carries identifiers the model may later cite. The
reconciler tool shows the pattern most clearly. The model asks for \code{compute\_reconciliation} with
optional bucket identifiers. Deterministic code computes the required, supported, gap and pending quantities
with exact decimals, and the model receives a structured summary. The model is never asked to calculate. If it
requests an operation outside the whitelist, the run stops with an expert-review outcome.

\subsection{Agentic RAG}
Retrieval is a tool that the agent chooses to use. The search tool returns scored passages with a grade for
authority, temporal currency and relevance. A lookup tool resolves an exact legal path to a provision
revision. When a grade is weak, the agent may reformulate the query at most twice for a finding. When
authority or date is still in doubt, the finding is escalated, and the system states the doubt in the report
and does not rely on the passage. Passages that are only leads cannot support a definitive statement.

\subsection{Planning and Reasoning}
The planning call returns a plan, with objective, hypotheses, proposed checks, expected evidence and
stopping criteria, and either an action or a stop reason. The assessment call returns a verdict
(supports, refutes, weak, no progress or unknown), a disposition and an optional revised plan. A plan record
is stored before each action and an assessment record after it, and each turn manifest says whether the plan
changed. The loop is therefore visible in the trace without storing private reasoning text.

\subsection{Observability}
Every model call stores the full request body and the full response, the token counts, the latency, the
model revision and the weights digest. Every tool call stores its intent and result. The user interface shows
a timeline of turns, the tool calls with their arguments and the model calls with their prompts, and the
HTTP API streams events over SSE. The same records feed the evaluation tables, so the numbers in
Chapter~\ref{ch:eval} are computed from the trace and are not entered by hand.

\subsection{Evaluation}
The evaluation has an independent reference for every case, a scorer in a separate process, an exact metric for
the primary outcome, and a human rubric for what code cannot judge. Section~\ref{sec:metrics} gives the
definitions.

\subsection{Bonus Concepts}
\textbf{Human-in-the-loop and guardrails.} The reviewer decision is bound to the exact report and input digests,
so a changed report needs a new packet. The guardrails include the tool whitelist, the finding gate,
the budget limits and the escape of markup characters in dossier text.
\textbf{Durable workflow.} The graph commits one turn at a time to an append-only store with a
compare-and-swap version. Section~\ref{sec:crash} reports a test in which the coordinator and its child were
killed mid-trial and the trial was retried.
\textbf{Experimentation.} The comparison of fixed and adaptive investigation is the planned experiment of
Chapter~\ref{ch:exp}, and development-set prompt experiments are reported in Section~\ref{sec:promptdev}.
